\section{Introduction}
\IEEEPARstart{F}{raud} screening is a temporal decision problem. Transactions arrive in order, labels often arrive later, and illicit activity is rare. Relational context can be useful, but its availability depends on how a production graph is maintained at scoring time. Investigators then inspect only a small part of a ranked queue. These conditions make the test set, graph view, and decision rule part of the scientific question, not details to be fixed after a model has been chosen. They also make precision--recall analysis and capacity-aware measures more informative than an isolated AUROC number \cite{saito2015pr,bahnsen2015cost,siddiqui2018feedback}.

Public financial-graph datasets have enabled increasingly capable detectors. Elliptic introduced temporally indexed Bitcoin transactions \cite{weber2019aml}; DGraphFin supplied a much larger account graph \cite{huang2022dgraph}; and IBM AML-Data made controlled transaction regimes available at several scales \cite{altman2023ibmaml}. Yet a result on any of these datasets remains conditional. A transductive graph exposes a different information set from a graph built only from past transactions. Removing cross-period edges tests a different reliance on structure than removing edge attributes. A threshold chosen for F1 addresses a different decision from a fixed review budget. If a configuration exhausts memory, the experiment has established a feasibility boundary, not poor predictive performance.

This dependence is easy to lose in a leaderboard. Statements such as ``graphs help'' or ``model $A$ is best'' omit the temporal horizon, visibility rule, graph construction, selection procedure, metric, and resource envelope under which the comparison was observed. Model-selection bias and data leakage can then enter through a split, a preprocessing step, or repeated inspection of test results \cite{cawley2010selection,kaufman2012leakage,kapoor2023leakage}. The broader benchmarking literature has shown that task and protocol choices can change apparent progress \cite{dehghani2021lottery,poursafaei2022better,reuel2024betterbench}. Temporal graph benchmarks make several of these choices explicit \cite{huang2023tgb,gastinger2024tgb2,huang2024benchtemp}, while graph-anomaly benchmarks broaden datasets and methods \cite{tang2023gadbench,hua2026bag}. Fraud screening still requires their concerns to be joined with graph construction, extreme prevalence, analyst capacity, and failed compute cells.

FraudShiftBench treats an evaluation result as a statement about a \emph{deployment contract}
\(
\Pi=(\mathcal{T},\mathcal{V},\mathcal{C},\mathcal{S},\mathcal{B},\mathcal{R}),
\)
whose components specify time, graph visibility, construction, model selection, decision budget, and resources. A result becomes claimable only through an evidence unit that records the prediction task, concrete contract, model and hyperparameters, seeds, metrics, prediction manifest, resource record, and integrity state. We then define a typed claim and a support relation between claims and evidence. This separation matters: insufficient evidence yields a blocked status; a compute failure yields resource-blocked; and an empirical predicate contradicted by complete evidence is refuted. None of those outcomes is silently converted into another.

The evidence does not reduce to one verdict about graphs. On Elliptic, strict versus isolated graph visibility reverses the top AUPRC ranking among MLP, GCN, and GraphSAGE. The same intervention leaves the MLP scores unchanged, improves GCN on Elliptic, and reduces GraphSAGE AUPRC on both real datasets. On synthetic IBM AML-Data, a strong tree baseline leads all eight baseline AUPRC cells, while the wider model and construction grid yields frequent AUPRC--fixed-threshold-F1 disagreement. Original transaction attributes, graph scale, recency, and degree capping have effects that depend on the cell. These results contradict both a universal graph benefit and a universal graph penalty. The measured resource boundary is equally specific: GINE is strongest by AUPRC in the four feasible Small graph-grid cells, but no Medium GINE score exists under the recorded T4 envelope.

This paper makes four contributions.
\begin{enumerate}
    \item It formalizes a fraud-specific deployment contract over six independent evaluation axes. The contract distinguishes temporal prediction, structural access, model selection, analyst capacity, and feasibility without treating one protocol as a universal default.
    \item It defines a typed claim--evidence support relation and an executable validator. Scope expansion, missing predictions, incomplete seeds, integrity failure, and resource failure have different semantics; controlled mutations of locked artifacts test those transitions.
    \item It provides a ten-seed empirical study across Elliptic, DGraphFin, and four IBM AML-Data Small/Medium regimes. The analysis isolates protocol, architecture, construction, metric, prevalence, scale, and runtime effects instead of pooling unlike prediction tasks into one ranking.
    \item It releases an auditable benchmark artifact that connects manuscript numbers to predictions, locks, and deterministic generators. A saved-output GraphSafe-TTA analysis is retained as a bounded decision case study, with negative and non-significant comparisons reported alongside favorable ones.
\end{enumerate}

The results address six questions: how graph-visibility protocols change conclusions (RQ1), whether those effects depend on architecture and dataset (RQ2), how IBM graph construction and scale alter rankings (RQ3), when rank and decision metrics disagree (RQ4), which performance--resource boundaries are observed (RQ5), and whether the support validator prevents unsupported promotion while preserving reproducibility (RQ6).
